% !PRINTED 616
\weberpage{616}

It only remains to show how one can determine $\varphi_2$ in accordance with the requirement laid down here. Since two forms which differ from one another only by a numerical factor are equivalent, $\varphi_2$ is equivalent to
\[
c_2\, (t_1 + \omega_2 t_2),
\]
[\S.\,166, (10)], and if we take this form for $\varphi_2$, then $b_{1,1} = c_2$. But the number $\omega_2$ can always be chosen among those equivalent to it in such a way that $c_2$ becomes relatively prime to $a_{1,1}$, and then $a_{1,1}$ is also relatively prime to $\varphi_2$.

For if one replaces $\omega_2$ by
\[
\frac{p\, \omega_2 + q}{r\, \omega_2 + s}, \quad p s - q r = \pm 1,
\]
then by Bd.\ I, \S.\,122, (13), $c_2$ passes over into
\[
c_2' = -a_2 r^2 + b_2 r s + c_2 s^2
\]
and one can now choose the integral rational numbers $r, s$ in such a way that they are coprime to one another and that $c_2'$ is not divisible by any given primes whatever.

For if $n$ is such a prime, then one takes, when $n$ divides $a_2$ and $c_2$, and hence does not divide $b_2$, $r$ and $s$ not divisible by $n$; when $n$ does not divide $a_2$, $s$ divisible by $n$ and $r$ not divisible by $n$; and when $n$ does not divide $c_2$, $r$ divisible by $n$, $s$ not divisible by $n$. (If $a_2$ and $c_2$ are both not divisible by $n$, one has the choice between the two latter assumptions.) These requirements are compatible with one another for different arbitrarily given primes $n$, and $p, q$ can then be determined from the condition $p s - q r = \pm 1$.

We have derived these laws of the composition of the quadratic irrational numbers here only under the assumption that the discriminant is at the same time the fundamental number of a quadratic field, i.e. that $\Delta$ is a stem discriminant. But the same laws also hold, with slight modifications, in general $^{1)}$.

\noindent\rule{4cm}{0.4pt}

$^{1)}$ \gesperrt{Gauss}, Disquisitiones arithm., Art.\,234 ff. \gesperrt{Dirichlet-Dedekind}, Zahlentheorie. 4. Auflage, \S.\,187.

\begin{center}
\rule{0.2\linewidth}{0.4pt}
\end{center}
